\documentclass[11pt]{article}
\usepackage[margin=1in]{geometry}
\usepackage{amsmath,amssymb,amsthm}
\usepackage{graphicx}
\usepackage{hyperref}

\newtheorem{theorem}{Theorem}
\newtheorem{proposition}{Proposition}
\newtheorem{lemma}{Lemma}
\newtheorem{remark}{Remark}

\newcommand{\E}{\mathbb E}
\newcommand{\Pp}{\mathbb P}
\newcommand{\R}{\mathbb R}
\newcommand{\cB}{\mathcal B}
\newcommand{\norm}[1]{\left\lVert #1\right\rVert}
\newcommand{\Lq}{L^q(S)}

\title{A \(q>2\) Deterministic-Dominator Square-Function Subcase}
\author{Agent \texttt{agent\_lane\_19}}
\date{2026-06-16}

\begin{document}
\maketitle

\section*{Status}

This packet records a narrow \(q>2\) partial result for Remark 12.10 of
Yaroslavtsev's paper \emph{Local characteristics and tangency of
vector-valued martingales} \cite{Yaroslavtsev2019}.  The preceding packet
\[
\texttt{1907.11588\_Lq\_p1\_endpoint\_qle2\_square\_function}
\]
proves the \(L^q(S)\), \(p=1\) endpoint for \(1<q\le2\).  The full endpoint
for \(q>2\) remains open here, but the square-function bottleneck is proved
when the dominating square magnitudes are deterministic.

\section*{Source Problem}

The source is arXiv:1907.11588, PDF page 105, Remark 12.10.  The source states
that the general local-martingale case is open and reduces the missing issue
to an independent-increment comparison.

\begin{center}
\includegraphics[width=\textwidth]{figures/open_problem_crop.png}
\end{center}

The continuation states the aggregate-domination formulation and notes that
even the symmetric case was not known to the author.

\begin{center}

\includegraphics[width=\textwidth]{figures/open_problem_continuation_crop.png}
\end{center}

\section*{Main Square-Function Subcase}

Let \(q>2\) and put \(r=q/2>1\).  The endpoint square-function comparison can
be written at the positive \(L^r(S)\)-level: for \(X_i\in L^q(S)\), set
\(V_i=|X_i|^2\in L^r(S)_+\).  The desired quantity is
\[
  \E\norm{\sum_i V_i}_{L^r(S)}^{1/2}
  =
  \E\left\|
      \left(\sum_i |X_i|^2\right)^{1/2}
    \right\|_{L^q(S)}.
\]

\begin{theorem}[Deterministic square magnitudes]
\label{thm:det}
Let \(q>2\), \(r=q/2\), and let \(v_1,\ldots,v_m\in L^r(S)_+\) be
deterministic.  Let \(V_1,\ldots,V_n\) be independent \(L^r(S)_+\)-valued
random variables.  Suppose that their aggregate occurrence measure is
dominated by the deterministic square-magnitude measure:
\[
  \sum_{i=1}^n \Pp(V_i\in A)
  \le
  \sum_{j=1}^m \mathbf 1_A(v_j),
  \qquad A\in \cB(L^r(S)),\ 0\notin A.
\]
Then
\[
  \E\norm{\sum_{i=1}^n V_i}_{L^r(S)}^{1/2}
  \le
  C_q
  \norm{\sum_{j=1}^m v_j}_{L^r(S)}^{1/2}.
\]
\end{theorem}

Equivalently, if \(u_j\in L^q(S)\) are deterministic and the aggregate laws
of \(|X_i|^2\) are dominated by the atoms \(|u_j|^2\), then
\[
  \E\left\|
      \left(\sum_i |X_i|^2\right)^{1/2}
    \right\|_{L^q(S)}
  \le C_q
  \left\|
      \left(\sum_j |u_j|^2\right)^{1/2}
    \right\|_{L^q(S)}.
\]

\section*{Idea of the Proof}

The \(1<q\le2\) packet used a positive multivariate Laplace representation.
That representation fails for \(q>2\).  In the deterministic-dominator
subcase one can instead square first and use scalar Rosenthal pointwise in
the spatial variable.

The scalar inequality controls
\[
  \E\left(\sum_i Z_i\right)^r
\]
by the sum of a mean term and an \(r\)-moment term.  Deterministic aggregate
domination controls both terms pointwise:
\[
  \sum_i \E V_i \le \sum_j v_j,
  \qquad
  \sum_i \E V_i^r \le \sum_j v_j^r\le\left(\sum_j v_j\right)^r.
\]
This is exactly the extra structure missing for a general random dominating
family.

\section*{Proof}

\begin{lemma}[Scalar positive Rosenthal]
\label{lem:scalar}
For every \(r>1\) there is a constant \(C_r\) such that for all finite
independent nonnegative scalar random variables \(Z_1,\ldots,Z_n\),
\[
  \E\left(\sum_i Z_i\right)^r
  \le
  C_r\left[
      \left(\sum_i \E Z_i\right)^r
      +
      \sum_i \E Z_i^r
    \right].
\]
\end{lemma}

\begin{proof}
For \(1<r\le2\), this is the von Bahr--Esseen inequality applied to
\(Z_i-\E Z_i\), followed by the triangle inequality in \(L^r\).  For \(r>2\),
the usual scalar Rosenthal inequality gives
\[
  \E\left|\sum_i(Z_i-\E Z_i)\right|^r
  \le
  C_r\left[
    \left(\sum_i\E Z_i^2\right)^{r/2}
    +
    \sum_i\E Z_i^r
  \right].
\]
Let \(A=\sum_i\E Z_i\) and \(B=\sum_i\E Z_i^r\).  Interpolation and Holder's
inequality give
\[
  \sum_i\E Z_i^2
  \le
  A^{(r-2)/(r-1)}B^{1/(r-1)}.
\]
Hence
\[
  \left(\sum_i\E Z_i^2\right)^{r/2}\le A^r+B
\]
by Young's inequality.  Combining this with
\(\|\sum_iZ_i\|_{L^r}\le A+\|\sum_i(Z_i-\E Z_i)\|_{L^r}\) proves the lemma.
\end{proof}

\begin{proof}[Proof of Theorem \ref{thm:det}]
Let
\[
  A=\sum_{j=1}^m v_j\in L^r(S)_+.
\]
The aggregate domination hypothesis implies the following two lattice
inequalities:
\[
  \sum_i \E V_i \le A
  \quad\text{in }L^r(S),
  \qquad
  \sum_i \E V_i^r \le \sum_j v_j^r
  \quad\text{in }L^1(S).
\]
Indeed, both follow by testing against nonnegative simple functions and using
the aggregate domination for the Borel maps
\[
  V\mapsto \int_S h\,V\,d\sigma,
  \qquad
  V\mapsto \int_S h\,V^r\,d\sigma
\]
with \(h\ge0\).

For almost every \(s\in S\), apply Lemma \ref{lem:scalar} to the independent
nonnegative scalar variables \(V_i(s)\).  Then
\[
  \E\left(\sum_i V_i(s)\right)^r
  \le
  C_r\left[
      \left(\sum_i \E V_i(s)\right)^r
      +
      \sum_i \E V_i(s)^r
    \right].
\]
Integrating over \(S\) gives
\[
\begin{aligned}
  \E\norm{\sum_i V_i}_{L^r}^r
  &\le
  C_r\left[
      \int_S\left(\sum_i \E V_i\right)^r\,d\sigma
      +
      \int_S \sum_i \E V_i^r\,d\sigma
    \right]  \\
  &\le
  C_r\left[
      \int_S A^r\,d\sigma
      +
      \int_S \sum_j v_j^r\,d\sigma
    \right]
  \le
  2C_r\norm{A}_{L^r}^r,
\end{aligned}
\]
because \(\sum_j v_j^r\le(\sum_j v_j)^r=A^r\).

Finally, by Jensen's inequality for the concave map \(t\mapsto t^{1/(2r)}\),
\[
  \E\norm{\sum_i V_i}_{L^r}^{1/2}
  =
  \E\left(\norm{\sum_i V_i}_{L^r}^{r}\right)^{1/(2r)}
  \le
  \left(\E\norm{\sum_i V_i}_{L^r}^{r}\right)^{1/(2r)}
  \le
  (2C_r)^{1/(2r)}\norm{A}_{L^r}^{1/2}.
\]
This is the asserted estimate.
\end{proof}

\section*{Discrete Martingale Corollary}

\begin{proposition}[Random signs with deterministic magnitudes]
\label{prop:signs}
Let \(q>2\).  Let \(u_1,\ldots,u_m\in L^q(S)\) be deterministic and let
\(\eta_j=\varepsilon_j u_j\), where \((\varepsilon_j)\) are independent
Rademacher signs.  Let \((\xi_i)\) be independent mean-zero \(L^q(S)\)-valued
random variables.  Assume that the aggregate square-magnitude laws satisfy
\[
  \sum_i \Pp(|\xi_i|^2\in A)
  \le
  \sum_j \mathbf 1_A(|u_j|^2),
  \qquad A\in\cB(L^{q/2}(S)),\ 0\notin A.
\]
Then
\[
   \E\max_{1\le k\le n}
      \left\|\sum_{i=1}^k \xi_i\right\|_{L^q(S)}
   \le
   C_q
   \E\max_{1\le k\le m}
      \left\|\sum_{j=1}^k \eta_j\right\|_{L^q(S)}.
\]
\end{proposition}

\begin{proof}
Theorem \ref{thm:det} gives
\[
  \E\left\|\left(\sum_i|\xi_i|^2\right)^{1/2}\right\|_{L^q}
  \le
  C_q
  \left\|\left(\sum_j|u_j|^2\right)^{1/2}\right\|_{L^q}.
\]
The right-hand side is, deterministically,
\[
  \E\left\|\left(\sum_j|\eta_j|^2\right)^{1/2}\right\|_{L^q}.
\]
Yaroslavtsev's Banach-function-space BDG theorem
\cite[Theorem 4.32 and Proposition 4.33]{Yaroslavtsev2018} compares the
\(L^1\) maximal norm of each martingale difference sequence with its lattice
square function.  Applying it to \((\xi_i)\) and \((\eta_j)\) proves the
claim.
\end{proof}

\section*{What This Does Not Yet Prove}

The full \(q>2\) endpoint would require the same square-function comparison
when the dominating family has random square magnitudes \(W_j=|Y_j|^2\).  The
argument above gives, for a dominated family \(V_i\),
\[
  \E\norm{\sum_i V_i}_{L^r}^{1/2}
  \le
  C_r
  \left[
      \int_S
      \left(\sum_j \E W_j\right)^r\,d\sigma
      +
      \int_S\sum_j \E W_j^r\,d\sigma
  \right]^{1/(2r)}.
\]
For deterministic \(W_j=v_j\), this is controlled by
\(\norm{\sum_j v_j}_{L^r}^{1/2}\).  For random \(W_j\), however, the target is
\[
  \E\norm{\sum_j W_j}_{L^r}^{1/2},
\]
and rare large jumps show that this small moment need not control the
displayed \(r\)-moment Rosenthal functional.  Thus this packet isolates a
specific anti-concentration or de-Poissonization step as the remaining
obstruction for this route.

\section*{Stress Tests and Failed Tickets}

During this push I also tested finite-dimensional regrouping patterns in
\(\ell_q^d\), including the \(c_0\)-style occupancy construction with fixed
finite \(q\).  The simple occupancy ratios stayed bounded, as expected:
for deterministic disjoint atoms, multinomial counts have uniformly bounded
\(r\)-moments for each fixed \(r=q/2\).  Random sparse/correlated atom searches
for \(r=1.5\) and \(r=2\) found only ratios around \(1.06\) and \(1.13\), not
dimension growth.  These computations are not proof, but they support the
view that the \(q>2\) obstruction is an anti-concentration proof issue rather
than the elementary \(c_0\) occupancy obstruction.

Two proof tickets remain plausible:
\begin{enumerate}
\item a positive-vector de-Poissonization theorem comparing
      \(\E\norm{\sum_i V_i}_{r}^{1/2}\) with the compound-Poisson model
      determined only by the aggregate measure;
\item a lower bound showing that every independent dominating family controls
      the correct truncated Rosenthal functional, not the too-large
      \(r\)-moment functional displayed above.
\end{enumerate}

\section*{Novelty Notes and Human Review}

The run indexes contain the \(1<p\) \(L^q\) packet, the \(1<q\le2\) endpoint
packet, Hilbert packets, the \(c_0\) counterexample, and the conditional
endpoint reduction.  No packet in the run records this deterministic
dominating-magnitude \(q>2\) subcase.

Human review should first check the lattice consequences of aggregate
domination:
\[
  \sum_i \E V_i\le\sum_jv_j,\qquad
  \sum_i \E V_i^r\le\sum_jv_j^r.
\]
After that, the rest is the scalar Rosenthal inequality and Jensen's
inequality.

\begin{thebibliography}{9}

\bibitem{Yaroslavtsev2019}
I. S. Yaroslavtsev.
\newblock Local characteristics and tangency of vector-valued martingales.
\newblock arXiv:1907.11588, 2019.

\bibitem{Yaroslavtsev2018}
I. S. Yaroslavtsev.
\newblock Burkholder--Davis--Gundy inequalities in UMD Banach spaces.
\newblock arXiv:1807.05573, 2018.

\bibitem{DirksenYaroslavtsev2017}
S. Dirksen and I. S. Yaroslavtsev.
\newblock \(L^q\)-valued Burkholder--Rosenthal inequalities.
\newblock arXiv:1707.00109, 2017.

\bibitem{Rosenthal1970}
H. P. Rosenthal.
\newblock On the subspaces of \(L_p\), \(p>2\), spanned by sequences of
  independent random variables.
\newblock Israel Journal of Mathematics 8 (1970), 273--303.

\bibitem{MontgomerySmithPruss}
S. J. Montgomery-Smith and A. R. Pruss.
\newblock A comparison inequality for sums of independent random variables.
\newblock arXiv:math/9811124, 1998.

\end{thebibliography}

\end{document}
